Mapa průměrných mezd v~sektoru obchodu a~oprav motorových vozidel (\ac{NACE} G) v~EU27 umísťuje ČR mezi státy s~nejnižší mzdovou úrovní v~tomto odvětví --- paradoxně v~situaci, kdy je tento sektor největším zaměstnavatelem v~absolutním počtu osob.
Maloobchod a~velkoobchod jsou v~ČR charakteristické téměř úplnou absencí \ac{KS} s~plošnou platností: velké maloobchodní řetězce (většinou zahraniční vlastnictví) vyjednávají individuálně nebo mají podnikové \ac{KS} s~minimálními tarify, zatímco tisíce menších zaměstnavatelů nejsou vázány žádnou kolektivní smlouvou.
Tato situace vytváří přímou soutěž o~pracovní sílu na bázi minimálních mzdových standardů: firmy, které by chtěly platit více, jsou vystaveny konkurenci firem, které platí méně --- race to the bottom bez koordinačního mechanismu.
Odvětvová \ac{KS} s~erga omnes v~sektoru G by vytvořila jednotný mzdový standard pro všechny maloobchodní zaměstnavatele --- eliminovala by mzdový dumping a~umožnila by lepším zaměstnavatelům soutěžit o~pracovníky na bázi nadstandardu, nikoliv race to the bottom.

Mapa průměrných mezd v~sektoru obchodu a~oprav motorových vozidel (\ac{NACE} G) v~EU27 umísťuje ČR mezi státy s~nejnižší mzdovou úrovní v~tomto odvětví --- paradoxně v~situaci, kdy je tento sektor největším zaměstnavatelem v~absolutním počtu osob.
Maloobchod a~velkoobchod jsou v~ČR charakteristické téměř úplnou absencí \ac{KS} s~plošnou platností: velké maloobchodní řetězce (většinou zahraniční vlastnictví) vyjednávají individuálně nebo mají podnikové \ac{KS} s~minimálními tarify, zatímco tisíce menších zaměstnavatelů nejsou vázány žádnou kolektivní smlouvou.
Tato situace vytváří přímou soutěž o~pracovní sílu na bázi minimálních mzdových standardů: firmy, které by chtěly platit více, jsou vystaveny konkurenci firem, které platí méně --- race to the bottom bez koordinačního mechanismu.
Odvětvová \ac{KS} s~erga omnes v~sektoru G by vytvořila jednotný mzdový standard pro všechny maloobchodní zaměstnavatele --- eliminovala by mzdový dumping a~umožnila by lepším zaměstnavatelům soutěžit o~pracovníky na bázi nadstandardu, nikoliv race to the bottom.

Mapa průměrných mezd v~sektoru obchodu a~oprav motorových vozidel (\ac{NACE} G) v~EU27 umísťuje ČR mezi státy s~nejnižší mzdovou úrovní v~tomto odvětví --- paradoxně v~situaci, kdy je tento sektor největším zaměstnavatelem v~absolutním počtu osob.
Maloobchod a~velkoobchod jsou v~ČR charakteristické téměř úplnou absencí \ac{KS} s~plošnou platností: velké maloobchodní řetězce (většinou zahraniční vlastnictví) vyjednávají individuálně nebo mají podnikové \ac{KS} s~minimálními tarify, zatímco tisíce menších zaměstnavatelů nejsou vázány žádnou kolektivní smlouvou.
Tato situace vytváří přímou soutěž o~pracovní sílu na bázi minimálních mzdových standardů: firmy, které by chtěly platit více, jsou vystaveny konkurenci firem, které platí méně --- race to the bottom bez koordinačního mechanismu.
Odvětvová \ac{KS} s~erga omnes v~sektoru G by vytvořila jednotný mzdový standard pro všechny maloobchodní zaměstnavatele --- eliminovala by mzdový dumping a~umožnila by lepším zaměstnavatelům soutěžit o~pracovníky na bázi nadstandardu, nikoliv race to the bottom.

Mapa průměrných mezd v~sektoru obchodu a~oprav motorových vozidel (\ac{NACE} G) v~EU27 umísťuje ČR mezi státy s~nejnižší mzdovou úrovní v~tomto odvětví --- paradoxně v~situaci, kdy je tento sektor největším zaměstnavatelem v~absolutním počtu osob.
Maloobchod a~velkoobchod jsou v~ČR charakteristické téměř úplnou absencí \ac{KS} s~plošnou platností: velké maloobchodní řetězce (většinou zahraniční vlastnictví) vyjednávají individuálně nebo mají podnikové \ac{KS} s~minimálními tarify, zatímco tisíce menších zaměstnavatelů nejsou vázány žádnou kolektivní smlouvou.
Tato situace vytváří přímou soutěž o~pracovní sílu na bázi minimálních mzdových standardů: firmy, které by chtěly platit více, jsou vystaveny konkurenci firem, které platí méně --- race to the bottom bez koordinačního mechanismu.
Odvětvová \ac{KS} s~erga omnes v~sektoru G by vytvořila jednotný mzdový standard pro všechny maloobchodní zaměstnavatele --- eliminovala by mzdový dumping a~umožnila by lepším zaměstnavatelům soutěžit o~pracovníky na bázi nadstandardu, nikoliv race to the bottom.

\input{texparts/python/sector_wages_map_G}



